\section{The small-flow bypass}
\label{sec:bypass}

A \emph{small flow} is one whose size is at most a threshold, 10\,KB in most of our experiments. The sending host marks every packet of a small flow, and its ACKs, with a DSCP code point (AF11), so routers can tell small-flow packets apart without inspecting anything else. Every router forwards a marked packet to a neighbor one hop closer to the destination, hashing the flow over all such neighbors (ECMP over shortest paths), so the packet stays on shortest paths end to end. Between neighboring routers that is the direct link.

Large flows keep using Spraypoint unchanged. The bypass adds no links and does not change Spraypoint's rings or tables; it is a separate lane on the same wires, as a separate VRF or traffic class would be on a real switch.

\paragraph{Implementation.} The simulator extends Spraypoint's per-router routing with one decision (\file{src/tcp-bypass/bypass-routing.cc}): for a marked packet, return the shortest-path next hop (\file{src/tcp-bypass/shortest-paths.cc}, which precomputes distances by breadth-first search and hashes over the closer neighbors); otherwise defer to Spraypoint. Flows, TCP, and measurements are shared with the baseline, so runs with and without the bypass carry identical traffic.

\subsection{Overhead on a switch (H4)}
\label{sec:h4}

Spraypoint selects its spray or point table by ingress port. The bypass adds a DSCP classification and a shortest-path ECMP table with up to $n-1$ destination entries per router. Its ECMP groups compete for the same resources that constrain Spraypoint's pointing groups (§7.4 of~\cite{rng}). A switch with sufficient table and action capacity could make this decision in one pipeline pass. The number of physical stages and the processing delay depend on the target and its resource occupancy. The simulator models queueing and forwarding paths without switch processing delay, so a target compilation and hardware measurements are still needed to evaluate H4. See \file{docs/OVERHEAD.md} for the pipeline and state analysis.
